\section{Reproducibility and evidence boundary}\label{sec:reproducibility}

The paper has six replay boundaries.
\begin{center}
\small
\begin{tabular}{@{}>{\raggedright\arraybackslash}p{0.25\textwidth}
>{\raggedright\arraybackslash}p{0.46\textwidth}
>{\raggedright\arraybackslash}p{0.20\textwidth}@{}}
\toprule
Object & Quick repository verifier & Role \\
\midrule
Global \(q,F_2\), strict \(\PF_3\) &
\path{scripts/verify_pf3_reduction.py} & algebra, tail, and cross-check audit \\
Global \(\Cfour>0\) &
\path{scripts/verify_c4_confluent.py} & algebra, tail, and cross-check audit \\
Quotient reduction &
\path{scripts/verify_pf4_wronskian_reduction.py} & exact generic algebra \\
Transport completion &
\path{scripts/verify_pf4_transport_kernel.py} & exact algebra and checks \\
Origin order-five obstruction &
\path{scripts/verify_pf5_origin.py} & exact rational \\
Continuous strict-\(\PF_4\) separator &
\texttt{scripts/verify\_continuous\_}\newline
\texttt{pf4\_separator.py} & exact and directed \\
\bottomrule
\end{tabular}
\end{center}

The two theorem-producing compact covers are separate from the quick wrappers:
\begin{center}
\small
\begin{tabular}{@{}>{\raggedright\arraybackslash}p{0.27\textwidth}
>{\raggedright\arraybackslash}p{0.62\textwidth}@{}}
\toprule
Cover & Full directed command \\
\midrule
\(q,F_2\), 7731 cells &
\path{work/2026-07-12-riemann-kernel-pf34-classification/}
\path{certify_pf3_curvature.py} \\
\(\Cfour\), 8050 cells &
\path{work/2026-07-13-riemann-kernel-pf4-verification/}
\path{certify_c4.py} \\
\bottomrule
\end{tabular}
\end{center}
Both use Arb at 192-bit precision and seven retained theta terms. The release
manifest records the complete expected stdout hashes. Their identifying
prefixes are
\[
 \texttt{a52a360cdc5e0}\quad(q,F_2),\qquad
 \texttt{bf2436a454ece2c4}\quad(\Cfour).
\]
The PF3 quick wrapper separately uses a 256-bit Arb jet cross-check; the C4
quick wrapper uses exact SymPy algebra and mpmath diagnostics. Thus wrapper
precision is not reported as base-cover precision.

The compact inequalities are directed Arb calculations at 192-bit precision;
the tail uses explicit analytic envelopes and rational coefficient bounds.
The origin order-five certificate is independent: it uses only standard-library
integer arithmetic with outward rounding on a rational grid, Machin's formula,
alternating Taylor sums, three retained theta modes, and one geometric tail
bound. The quotient, curvature, primitive, endpoint, density-ratio, and
confluent-limit identities are mathematical identities for generic smooth
functions. The paper displays these identities; the scripts independently
replay them.

A focused independent audit should proceed in this order:
\begin{enumerate}
 \item verify the analytic theta normalization in \eqref{eq:kernel-series};
 \item replay the directed \(q,F_2,\Cfour\) compact and tail certificates;
 \item check the sign \(\T\log A=M\) and the two weighted-mean splits in
       \secref{sec:quotient};
 \item expand the determinant and curvature identity in
       \appref{app:algebra};
 \item check the endpoint identity in \appref{app:endpoints};
 \item differentiate the density ratio in
       \eqref{eq:density-ratio-derivative};
 \item derive the parity factorization and divided-difference limit in
       \secref{sec:completion}, then replay the rational origin certificate;
 \item reconstruct the separator's central determinant; then verify the
       coefficient signs and Fourier discriminant in \appref{app:separator}.
\end{enumerate}
No numerical scan is used as evidence for global nonconfluent positivity of
either kernel.

The non-mutating entry point \path{scripts/replay_paper.py} checks all six
registered verifier outputs in quick mode. The option \texttt{--full} first
reruns both compact covers. Exact package versions are in
\path{requirements-paper.txt}.
